
\documentclass[12pt]{article}
\usepackage[utf8]{inputenc}
\usepackage[T1]{fontenc}
\usepackage{lmodern}
\usepackage{microtype}
\usepackage[hidelinks]{hyperref}
\usepackage[a4paper, margin=3cm]{geometry}
\usepackage{parskip}

\title{Chaitra}
\date{\today}

\begin{document}
\maketitle

\section*{Why this letter?}

After I decided to leave Instagram last autumn, I promised everyone I would share updates through some kind of newsletter.
Reflecting on what took me so long yielded the following reasons. First of all, it was not clear to me what exactly I should be writing. What nice cakes I ate lately? Or where I have been traveling to? Or rather how my meditations were going? Or maybe all of it? Secondly, who I should send it to. Like we tell different people different things when they ask: "How are you doing?". It was not clear to me if there was one congruent answer I could give to everyone without making it too vague - and without wasting anyone's time. Thirdly, there is obviously an underestimated overhead with writing such a "newsletter" that requires something that can sometimes be hard to gather in modern times - yes, you guessed it. Time. It took me significant amount of time to sort myself and my task management system to become more efficient with my time and also more intentional.

Anyways, excuses aside. Today I want to share what I want to achieve with this newsletter. Even though I had a great 6 months without social media, where I took up new projects, read more books than before, and didn't miss out on anything, it is important to me to stay in touch with you amazing people and engage in meaningful exchange and rapport. It is important to me to stress this point. Honestly, if I reflect: What is the greatest asset I have for achieving anything that may emancipate this world? Where do I get inspiration from? Where do I go if I need any feedback or advice? Who would correct me if I end up on a destructive path? The answer is: You! My desire is to offer to you whatever I am doing at the moment and to have it blessed by your attention, or by your comment, or criticism. I'm convinced, then only, I can achieve anything meaningful.

The idea is to send out a letter every full moon. That’s why this one is called ‘Chaitra’ - also referred to as ‘the pink moon’, as it coincides with the blooming of spring wildflowers.

\section*{What am I up to at the moment?}

Let’s begin with something easy. I am 28 years old, live in Amsterdam with my wife, Ambika and our dog, Ollie. I am employed as a Software Engineer for a company called DNV and work for a project that is concerned with building a software to support the grid restoration process (i.e. the process of restarting the power grid after an outage or a blackout). Even though I’m employed in the Netherlands (since beginning of the year), I mainly work with German colleagues, who are located in Dresden. That’s why I travel to Dresden roughly every 4-6 weeks and work from home the rest of the time. Sometimes I also travel to the customer site or to conferences and other events, which I generally enjoy doing.

I also play around with AI regularly and thought to invite Claude to help me out here a bit. I do write a daily journal and gave Claude access to it so it give you a summary.
So Claude, how have I been lately?

\begin{quote}
\itshape
Honestly, pretty good — better than you've been in a while.

The year started with one of those rare moments of real clarity that made you realise you'd been a bit lost without fully noticing. Since then things have felt more purposeful. Brazil in January and February was a big part of that: nearly a month on the road with the Nirmalights, performing across five cities, and coming back with a clearer sense of what you want to build and why.

Since then you've been conducting a Sahaja Yoga course, started classical Indian vocal training — humbling and great in equal measure — and you're currently deep in organising a concert: Pandit Abhay Sopori at Amstelkerk in Amsterdam on May 22. That last one has taken over most of your spare hours.

Overall: engaged, more alive than a year ago, and with enough on your plate that you're rarely bored.
\end{quote}

By the way - this letter is 100\% human written unless clearly stated like above ;).

\section*{Some resources}

\textbf{Books I enjoyed recently}

\begin{itemize}
  \item \textit{Sapiens} — Yuval Noah Harari
  \item \textit{Moral Ambition} — Rutger Bregman
  \item \textit{Atomic Habits} — James Clear
  \item \textit{Building a Second Brain} — Tiago Forte
\end{itemize}

\textbf{Links}

\begin{itemize}
  \item My GitHub: \href{https://github.com/fhjgch?tab=repositories}{github.com/fhjgch}
  \item My Readwise shortlist: \href{https://readwise.io/reader/view/9b467e3f-5554-4749-9c2a-be7ffd087aee}{readwise.io}
  \item Latest blog post on meditation: \href{https://fhjgch.github.io/scriptorium/how-to-get-into-meditation.html}{How to get into meditation}
\end{itemize}

\end{document}
